
\subsection{3.1 模型基本假设}
为构建贴合赛题计费规则、物理可行且可复现的微网时序优化模型，本文在不改变题目核心约束的前提下设置如下建模假设：

\begin{enumerate}
    \item \textbf{电网交易规则}：微网仅支持向上级电网购电，无反向售电能力。富余光伏可弃光且无惩罚；已签约的日前合同电量若未完全消纳，作为闲置弃购电量按原价计费，电网不予以补偿抵扣。
    \item \textbf{储能参数约束}：储能额定容量为 $12000\ \mathrm{kWh}$，有效运行区间为 $[1200,10800]\ \mathrm{kWh}$。单时段最大充、放电功率均为 $5000\ \mathrm{kW}$，充放电效率取 $\eta_c=\eta_d=0.9$。后续通过效率参数替换完成敏感性分析。
    \item \textbf{调度行为约束}：同一调度时段不存在物理意义上的同时充放电行为。求解过程中出现的极小数值共生解将通过次级目标筛选剔除。
    \item \textbf{时序仿真规则}：问题一采用典型日闭环调度，约束日初与日末 SOC 相等；问题二至四采用全年连续仿真，取消每日电量闭环约束。引入终端价值罚项抑制有限时域优化短视性，仅用于优化引导，不参与费用统计。
    \item \textbf{非前视信息约束（核心）}：全文严格遵循非前视决策准则，所有预测、购电计划制定、合同调账仅可使用决策时刻之前的已知信息，禁止使用未来真实源荷、电价数据。
    \item \textbf{成本核算边界}：题目未给出设备损耗、潮流约束及老化成本，模型仅统计合同购电费用与应急购电费用，储能吞吐作为次级优化指标以提升调度平滑性。
\end{enumerate}

核心交易规则、非前视信息集为刚性边界假设；储能参数、数值处理规则为可检验建模设定，可通过敏感性分析验证稳健性。

\subsection{3.2 统一时间口径与量纲转换}
本文采用题目标准 10 分钟调度时序，单时段时长 $\Delta t=\displaystyle\frac{1}{6}\ \mathrm{h}$，通过 $E=P\cdot\Delta t$ 统一将功率量纲转换为电量，保证储能递推、能量平衡与费用计算量纲一致。

数据时刻标签为时段结束时刻，第 $t$ 行数据对应区间 $[(t-1)\times10,\ t\times10)\ \mathrm{min}$，实现全日时序无缝衔接。针对官方输出模板存在的时序偏移问题，本文仿真计算、SOC 更新均基于原始时序逻辑执行，并设置对照实验完成稳健性检验。

\subsection{3.3 基础功率平衡物理模型}
结合微网源储荷耦合特性，构建适用于全场景的时序优化模型。定义 $L_t$ 为时段负荷电量，$V_t$ 为可用光伏电量，$g_t$ 为生效合同购电量，$w_t$ 为闲置弃购电量，$v_t$ 为实际光伏消纳量，$x_t,y_t$ 为储能充、放电电量，$u_t$ 为应急购电量，$S_t$ 为时段初始储能电量。

时段能量平衡方程：
\[
g_t - w_t + v_t + y_t + u_t = L_t + x_t
\]

储能荷电状态递推关系：
\[
S_{t+1}=S_t+\eta_c x_t-\dfrac{y_t}{\eta_d},\quad \eta_c=\eta_d=0.9
\]

系统约束条件：
\[
\begin{cases}
0\le v_t \le V_t,\quad 0\le w_t \le g_t,\quad u_t \ge 0\\
0\le x_t,y_t \le 833.33\\
1200 \le S_t \le 10800
\end{cases}
\]

全年仿真采用连续 SOC 跨日传递机制，以前一日末电量作为次日初值。本文以1月为模型预热期，2—12月为正式统计周期，规避每日重置 SOC 导致的仿真失真。

\subsection{3.4 模型符号说明}
\begin{table}[htbp]
\centering
\caption{模型符号说明}
\begin{tabular}{ccc}
\hline
符号 & 单位 & 含义说明 \\
\hline
$L_t,V_t$ & kWh & 时段负荷电量、可用光伏电量上限 \\
$g_t,w_t$ & kWh & 生效合同购电量、闲置弃购电量 \\
$v_t,x_t,y_t,u_t$ & kWh & 光伏消纳量、储能充放电量、应急购电量 \\
$S_t$ & kWh & 时段初始储能荷电状态 \\
$c_t$ & 元/kWh & 分时电价 \\
$p_t,q_t$ & kWh & 初始日前合同电量、滚动调账后合同电量 \\
$\mathcal{F}_{d,h}$ & $-$ & 决策时刻可获取信息集 \\
\hline
\end{tabular}
\label{tab:symbol}
\end{table}

本文引入闲置弃购变量 $w_t$，严格区分合同采购电量与实际入网电量，精准匹配“按量计费、富余不退”的题目规则，保证能量平衡与经济核算逻辑自洽。